\section{An integral marking by smooth singular cycles}
\label{sec:integral-marking}

The reference fibre in Theorem~\ref{thm:evaluated-periods} is a
smooth complete intersection with algebraic coefficients. We now
construct an integral marking on that fibre, together with smooth
representatives suitable for integration. The construction uses the
global fibre, but does not express any local vanishing sphere in the
resulting basis.

\begin{theorem}\label{thm:smooth-marking}
Let $X\subset\PP^3\times\PP^3$ be a smooth complete intersection
of bidegrees $(1,1)$, $(1,1)$ and $(2,2)$ with represented algebraic complex
coefficients. There is a terminating procedure producing an integral
symplectic basis of $H_3(X;\ZZ)$ and, for each basis element, a finite
integral sum of smooth singular three-simplices in $X$.
The face cancellations are exact. Each simplex map has a specified
neighbourhood of the closed simplex on which it is smooth; its values
and derivatives through order two are evaluated with known finite bounds.
\end{theorem}

The real and imaginary parts of an algebraic coefficient are supplied
with their polynomial and isolating data. Exact real algebraic decisions
below apply to this representation. They are not decisions about
arbitrary computable real numbers.

\subsection{A compact real algebraic presentation}

Embed each factor $\PP^3$ in the real vector space of Hermitian
matrices by its rank-one orthogonal projector. It is described by
$P^*=P$, $P^2=P$, and $\tr P=1$. In the product embedding, the
squared norm of a homogeneous equation on unit representatives is
a polynomial in the two projector matrices: expand the squared
modulus and replace $z_i\overline z_j$ by $P_{ij}$ in each factor.
The three squared equations therefore give a compact real algebraic
presentation of $X$.

These squared equations need not have independent differentials.
For tangent spaces we retain the original complex equations in
projective affine charts. Their real derivative has rank six on
$X$. Its kernel, carried into the projector embedding by the
algebraic chart map, describes $TX$ effectively on a finite cover.

The compatible semialgebraic triangulation theorem gives a finite
simplicial complex $K$ and a semialgebraic homeomorphism
$h:|K|\longrightarrow X$, with algebraic descriptions of its graphs
on simplices \cite[Theorem~5.43 and Remark~11.46(b)]{BPR}.
Here and below the locators for this reference are those of its
2016 posted revision. We use the classical construction of a
map into $X$, rather than a simplicial replacement known only
to have the same homology. No differentiability of $h$ is assumed.

\subsection{An effective normal tube}

We give the required retraction construction for a compact smooth
real algebraic submanifold $M\subset\RR^N$ with effective local
regular equations. Its tangent descriptions give a semialgebraic
normal bundle
\[
 NM=\{(x,v):x\in M,\ v\perp T_xM\}.
\]
The map $E(x,v)=x+v$ has invertible differential along the zero
section. There exists $r_0>0$ such that it is locally invertible and
injective on $\|v\|<r_0$. For injectivity, otherwise equal-image
pairs with normal lengths tending to zero would, after passing to
a subsequence, have the same limiting basepoint in the compact
manifold. Both pairs would then lie in one inverse neighbourhood
of its zero normal vector, a contradiction. Compactness also makes
local invertibility uniform near the zero section.

For a positive rational $r$, the assertions that $E$ is injective
on $NM_{<r}$ and that its differential is invertible there are
first-order conditions over the real algebraic numbers. Injectivity
uses two normal vectors and equality of their images. Differential
invertibility is checked in the finite semialgebraic normal-bundle
charts; their derivative graphs are obtained by implicit
differentiation. Real quantifier elimination decides these conditions.
Testing $r=2^{-k}$ therefore terminates.

The image of the accepted tube is exactly
$T_r=\{y:\dist(y,M)<r\}$. A nearest point exists by compactness,
and the difference from a nearest point is normal. Injectivity of
$E$ proves that this point is unique for $y\in T_r$ and also
for every image $E(x,v)$ with $\|v\|<r$. The first component of
$E^{-1}$ is consequently a smooth nearest-point retraction
\begin{equation}\label{eq:normal-retraction}
                       r_M:T_r\longrightarrow M.
\end{equation}
Its graph is semialgebraic. This proves the particular retraction
we need; general semialgebraic tubular-neighbourhood existence
results alone would not supply its radius and derivative bounds.

On the compact closed tube of radius $3r/4$, every fixed derivative
of $r_M$ is bounded. Its graph is effectively semialgebraic, either
by implicit differentiation in inverse charts or by the first-order
definition of derivatives followed by quantifier elimination.
Enumerating integer component bounds and deciding the universal
inequalities produces finite bounds of any prescribed order.

At an algebraic point in $T_r$, the unique value of $r_M$, and of
each derivative, is algebraic and can be isolated. At a named point
$y$ in the smaller tube of radius $r/2$, choose an algebraic
$q$ with $\|q-y\|<r/8$. The joining segment lies in the tube
of radius $3r/4$. A bound for one further derivative controls the
error made by replacing $y$ by $q$. Refining $q$ and the algebraic
isolation therefore evaluates $r_M(y)$ and its required derivatives.
No sign or equality test on $y$ is used.

\subsection{Face-compatible smoothing}

Order the vertices of $K$ once. The graph of $h$ on a standard
closed $d$-simplex has algebraic data, so values at algebraic points
are algebraic. It also has an effective uniform modulus of
continuity: for any rational $\epsilon>0$, search for a rational
$\delta>0$ such that
\begin{equation}\label{eq:triangulation-modulus}
 \|u-v\|\le\delta\quad\Longrightarrow\quad
             \|h(u)-h(v)\|\le\epsilon
             \qquad (u,v\in\Delta^d).
\end{equation}
This is a first-order algebraic condition on the graph. Uniform
continuity proves termination. Take minima over the finite complex,
and obtain a common finite bound $R_0\ge1$ for $\|h\|$ in the
same way.

For barycentric coordinates $\lambda_0,\ldots,\lambda_d$, use
one common degree $n\ge1$ and the Bernstein polynomial
\begin{equation}\label{eq:bernstein-map}
 B_nh(\lambda)=\sum_{|\alpha|=n}
                 \frac{n!}{\alpha!}\lambda^\alpha h(\alpha/n).
\end{equation}
Its coefficients are algebraic. On a face all terms involving an
omitted vertex vanish; the remaining terms give exactly the same
polynomial as the construction on that face. The restrictions of
the maps thus agree as parametrized maps, not just as sets.

We record a uniform approximation bound. In the coordinates
$(\lambda_1,\ldots,\lambda_d)$, let $A$ have the multinomial
distribution with parameters $n,\lambda$. Then $B_nh$ is the
expectation of $h(A/n)$, and
\[
 \mathbb E\left\|(A_1/n,\ldots,A_d/n)
                    -(\lambda_1,\ldots,\lambda_d)\right\|^2
                      \le d/(4n).
\]
Choose $\delta$ from \eqref{eq:triangulation-modulus} for
$\epsilon/2$. Splitting at distance $\delta$ and applying
Markov's inequality gives
\begin{equation}\label{eq:bernstein-bound}
       \|B_nh-h\|_\infty\le\epsilon/2
                           +R_0d/(2n\delta^2).
\end{equation}
One common integer $n\ge\max_d R_0d/(\epsilon\delta^2)$
therefore gives error at most $\epsilon$ on the entire complex.

Apply the normal-tube construction to $M=X$ and choose
$\epsilon<r/8$. Define, simplex by simplex,
\begin{equation}\label{eq:smooth-simplex-map}
                         g=r_X\circ B_nh.
\end{equation}
It has the same exact face identities. It is smooth on a known
neighbourhood of each closed simplex: bound the polynomial derivatives
on a larger compact box and choose a sufficiently small neighbourhood
whose image stays in $T_{r/2}$. If $L$ bounds its first derivative,
one may take a radius smaller than the box margin and
$3r/[8(1+L)]$. The chain rule gives the evaluated bounds
\begin{equation}\label{eq:smoothed-chain-derivatives}
 \begin{split}
 \|Dg\|&\le\|Dr_X\|\,\|D(B_nh)\|,\\
 \|D^2g\|&\le\|D^2r_X\|\,\|D(B_nh)\|^2
                       +\|Dr_X\|\,\|D^2(B_nh)\|.
 \end{split}
\end{equation}
The face-compatible homotopy
\[
 (u,t)\longmapsto r_X\bigl((1-t)h(u)+tB_nh(u)\bigr)
\]
is defined because its argument stays within $\epsilon$ of $h(u)$.
It identifies $g_*$ with $h_*$ on integral homology. We assert
neither injectivity nor immersion for $g$. It produces smooth
singular chains and does not require a $C^2$ triangulation theorem.

\subsection{Integral homology and orientation}

The integer boundary matrices of the ordered finite complex compute
integral homology and cohomology by Smith reduction, including chain
representatives. The smooth-CICY topology gives
$H_1(X;\ZZ)=0$, $H_2(X;\ZZ)=\ZZ^2$, and free
$H_3(X;\ZZ)$ of rank $134$. Thus its middle intersection form is
unimodular and alternating. In particular these computations retain
integral indices; rational linear algebra would not suffice.

It remains to select the complex orientation. Choose either primitive
generator $z_6$ of $H_6(K;\ZZ)$. For the fixed product Fubini--Study
form $\omega_K$, put $\beta=(\omega_K|_X)^3/3!$. Since $g$ is
homotopic to the homeomorphism $h$,
\[
                    \int_{g_\#z_6}\beta=\pm\Vol_{\omega_K}(X)\ne0.
\]
This sign is determined by integration with an absolute error bound.
If $C_1,C_2$ bound the first two derivatives of a six-simplex map
and $W_0,W_1$ bound a smooth ambient extension of $\beta$ and
its first derivative, its pullback integrand has Lipschitz bound
\begin{equation}\label{eq:orientation-quadrature}
                    W_1C_1^7+6W_0C_1^5C_2.
\end{equation}
Subdivision and pointwise evaluation consequently give arbitrarily
small absolute enclosures after summing with the absolute chain
coefficients. Since the true integral is nonzero, an enclosure
eventually excludes zero. Reverse $z_6$ if necessary. This determines
the complex orientation using only the known ambient positive form.

\subsection{Intersection form and symplectic basis}

Let $\xi_1,\ldots,\xi_{134}$ be an integral cohomology basis
computed on $K$. The simplicial cup product gives
\[
             Q_{ab}=\langle\xi_a\smile\xi_b,z_6\rangle.
\]
Use the left cap convention specified by
\begin{equation}\label{eq:cap-convention}
 \langle\eta,\xi\cap z_6\rangle
                     =\langle\xi\smile\eta,z_6\rangle.
\end{equation}
Then $c_a=\xi_a\cap z_6$ are integral three-cycles whose classes
are a basis by Poincar\'e duality. Their oriented homological
intersection matrix is $Q$ in this convention. The cup/cap identities
and integral duality are those of
\cite[Theorem~3.30 and Proposition~3.38]{Hatcher}; the explicit
order in \eqref{eq:cap-convention} fixes the sign in degree three.

Integral symplectic reduction finds $S\in\operatorname{GL}_{134}(\ZZ)$
with $S^TQS$ equal to the chosen standard symplectic matrix. To see
termination directly, a primitive vector has a primitive pairing row
by unimodularity, so B\'ezout's identity finds a partner pairing to
one. Their hyperbolic plane splits integrally; the orthogonal
projection uses integer pairing values. Compute a basis of the
complement by Smith reduction and repeat.

Apply these same integer combinations to the cycles $c_a$, and then
apply $g_\#$. Exact face identities preserve every chain boundary
cancellation. Equations \eqref{eq:normal-retraction}--\eqref{eq:smoothed-chain-derivatives}
give the asserted domains, evaluations and bounds for their smooth
simplex maps. This proves Theorem~\ref{thm:smooth-marking}.

The resulting marking is constructed from the global fibre equations.
The period observations will identify the local vanishing spheres
in this marking.
